\ExerciseSection

\begin{enumerate}
\def\labelenumi{\arabic{enumi})}
\item
  Let \(f\) be the real-valued function on \(\mathbf{R}\) which is equal to \(0\) for \(x \leqslant 0\) , to \(x\) for \(0 \leqslant x \leqslant 1\) , and to \(1\) for \(x \geqslant 1\) . Show that the sequence of real-valued functions \(f_n(x) = f(nx - n^2)\) ( \(n \in \mathbf{N}\) ) is equicontinuous but not uniformly equicontinuous on \(\mathbf{R}\) , although it is formed of uniformly continuous functions.
\item
  Let X be a Hausdorff topological space in which every point admits a countable fundamental system of neighbourhoods, and let Y be a uniform space. Let H be a subset of \(\mathcal{C}(X;Y)\) with the property that, for each compact subset K of X, the set \(H|K\) of restrictions to K of mappings \(u\in H\) is an equicontinuous subset of \(\mathcal{C}(K;Y)\) . Show that H is equicontinuous (argue by contradiction).
\item
  Let X, Y, Z be three metric spaces and let H be a subset of \(\mathcal{C}(X \times Y; Z)\) . Suppose that, for each \(x_{0} \in X\) , the set of mappings \(u(x_{0}, .)\) , as u runs through H, is an equicontinuous subset of \(\mathcal{C}(Y; Z)\) , and that, for each \(y_{0} \in Y\) , the set of mappings \(u(., y_{0})\) , as u runs through
\end{enumerate}

H, is an equicontinuous subset of \(\mathcal{C}\) (X; Z). Show that if X is complete, then for each \(b \in \mathbf{Y}\) there is a set \(S_b\) in X whose complement is meagre in X and which is such that, for each \(a \in \mathbf{S}_b\) , the set H is equicontinuous at \((a, b)\) . (Apply Chapter IX, § 5, Exercise 23, using Chapter X, § 2, Proposition 1, Corollary 1.)

\begin{enumerate}
\def\labelenumi{\arabic{enumi})}
\setcounter{enumi}{3}
\tightlist
\item
  Let X, Y, Z be three uniform spaces.
\end{enumerate}

\begin{enumerate}
\def\labelenumi{\alph{enumi})}
\item
  Let H be a uniformly equicontinuous subset of \(\mathcal{C}(\mathrm{Y};\mathrm{Z})\) . If H, \(\mathcal{C}(\mathrm{X};\mathrm{Y})\) and \(\mathcal{C}(\mathrm{X};\mathrm{Z})\) are endowed with the uniformity of uniform convergence, show that the mapping \((u,v)\to u\circ v\) of \(\mathrm{H}\times \mathcal{C}(\mathrm{X};\mathrm{Y})\) into \(\mathcal{C}(\mathrm{X};\mathrm{Z})\) is uniformly continuous.
\item
  Let K be a uniformly equicontinuous subset of \(\mathcal{C}(\mathbf{X};\mathbf{Y})\) , and let L be a uniformly equicontinuous subset of \(\mathcal{C}(\mathbf{Y};\mathbf{Z})\) . Show that the set of mappings \(v \circ u\) , where u runs through K and v runs through L, is a uniformly equicontinuous subset of \(\mathcal{C}(\mathbf{X};\mathbf{Z})\) .
\end{enumerate}

\begin{enumerate}
\def\labelenumi{\arabic{enumi})}
\setcounter{enumi}{4}
\item
  Let X be a topological space and let Y be a normed space over a non-discrete valued division ring K. Let H be a subset of \(\mathcal{F}(X;Y)\) , equicontinuous at a point \(x_{0}\in X\) , and let k\textgreater0 be a real number. Then the set \(H_{k}\) of linear combinations \(\sum_{i}c_{i}u_{i}\) of functions \(u_{i}\in H\) , such that \(\sum_{i}|c_{i}|\leqslant k\) , is equicontinuous at the point \(x_{0}\) .
\item
  Let X be a topological space, Y a uniform space, H an equicontinuous subset of \(\mathcal{C}(\mathbf{X};\mathbf{Y})\) , and \(\Phi\) a filter on H. Show that the set of all \(x\in X\) such that \(\Phi(x)\) is a Cauchy filter base on Y is closed in X.
\item
  Let X be a topological space, Y a complete Hausdorff uniform space, and let \(\varphi\) be a uniformly continuous homeomorphism of Y onto an open subset \(\varphi(Y)\) of a Hausdorff uniform space \(Y'\) . Let H be an equicontinuous subset of \(\mathcal{C}(X; Y)\) , and let \(H'\) be the set of mappings \(\varphi \circ u\) , where \(u \in H\) . If \(v\) lies in the closure of \(H'\) in \(\mathcal{F}_s(X; Y')\) , show that \(\overline{v}^{1}(\varphi(Y))\) is both open and closed in X. In particular, if X is connected, \(v(X)\) is contained in \(\varphi(Y)\) or in a connected component of \(\mathbb{C}\varphi(Y)\) . (Observe that \(v\) is continuous, and use Exercise 6.)
\item
  Let \(\mathbf{H}\) be an equicontinuous set of mappings of a topological space \(\mathbf{X}\) into \(\mathbf{R}\) .
\end{enumerate}

\begin{enumerate}
\def\labelenumi{\alph{enumi})}
\item
  Show that the set of upper (resp. lower) envelopes of finite subsets of H is equicontinuous.
\item
  Let v be a mapping of X into \(\overline{R}\) which lies in the closure of H with respect to the topology of pointwise convergence. Show that v is continuous on \(\mathbf{X}\) and that the sets \(\overline{v} (+\infty)\) and \(\overline{v} (-\infty)\) are both open and closed in \(\mathbf{X}\) (use Exercise 7).
\item
  Deduce from a) and b) that the upper (resp. lower) envelope w (resp. v) of H is continuous on X and that \(\overline{w}^{1}(+\infty)\) {[}resp. \(\overline{v}^{1}(-\infty)\) {]} is both open and closed in X.
\item
  Suppose that X is connected and that there is a point \(x_{0} \in X\) such that \(\mathbf{H}(x_{0})\) is bounded above in R. Show that, for every compact subset K of X, the set of restrictions to K of functions of H is uniformly bounded above in K {[}use c){]}.
\end{enumerate}

\begin{enumerate}
\def\labelenumi{\arabic{enumi})}
\setcounter{enumi}{8}
\item
  Let X be a set, Y a uniform space and S a covering of X. A subset H of the uniform space \(\mathcal{F}_{\mathfrak{S}}(\mathbf{X};\mathbf{Y})\) is precompact if and only if (i) for each \(x\in X\) , \(\mathbf{H}(x)\) is precompact in Y; (ii) the uniformity of S-convergence and that of pointwise convergence coincide on H. (Apply Theorem 2 by considering X as a discrete space and reducing to the case where Y is Hausdorff and complete.)
\item
  \begin{enumerate}
  \def\labelenumii{\alph{enumii})}
  \tightlist
  \item
    Let X be a compact space, let \((f_{n})\) be a sequence of continuous mappings \(X \rightarrow X\) , which converges pointwise but not uniformly in X to a continuous function (§ 1, no. 6, Remark 2). Show that the set of mappings \(f_{n}\) is relatively compact in \(\mathcal{C}_{s}(X; X)\) but not in
  \end{enumerate}
\end{enumerate}

\[
\mathcal {C} _ {u} (\mathrm{X}; \mathrm{X}) = \mathcal {C} _ {c} (\mathrm{X}; \mathrm{X}).
\]

\begin{enumerate}
\def\labelenumi{\alph{enumi})}
\setcounter{enumi}{1}
\tightlist
\item
  Let I be the interval \([-1, +1]\) of R and for each integer n \textgreater{} 0 let \(u_{n}(x) = \sin \sqrt{x + 4n^{2}\pi^{2}}\) for all \(x \geqslant 0\) . Show that the set H of functions \(u_{n}\) is an equicontinuous subset of \(\mathcal{C}(\mathbf{R}_{+}; \mathbf{I})\) and is relatively compact in \(\mathcal{C}_{c}(\mathbf{R}_{+}; \mathbf{I})\) but not in \(\mathcal{C}_{u}(\mathbf{R}_{+}; \mathbf{I})\) . (Note that the sequence \((u_{n})\) converges pointwise to 0.)
\end{enumerate}

\begin{enumerate}
\def\labelenumi{\arabic{enumi})}
\setcounter{enumi}{10}
\item
  Let X be a completely regular space and let S be a set of subsets of X whose interiors cover X. Show that the space \(\mathcal{C}_{s}(X;R)\) is not locally compact.
\item
  Let X be a topological space, Y a uniform space, H an equicontinuous subset of C(X; Y) and V a symmetric entourage of Y. Given a point \(x \in X\) , show that the set of points \(x' \in X\) for which there is an integer n (depending on \(x'\) ) such that \(\mathrm{H}(x') \subset \mathrm{V}^n(\mathrm{H}(x))\) is both open and closed in X. Deduce that, if K is any compact connected subset of X and if \(x_0\) is any point of K, there is an integer n \textgreater{} 0 such that \(\mathrm{H}(\mathrm{K}) \subset \mathrm{V}^n(\mathrm{H}(x_0))\) .
\end{enumerate}

¶ 13) Let X be a topological space, let Y be a locally compact uniform space and let H be an equicontinuous subset of C(X;Y).

\begin{enumerate}
\def\labelenumi{\alph{enumi})}
\item
  Let A be the set of all \(x \in X\) such that \(H(x)\) is relatively compact in Y. Show that A is open in X (cf.~Chapter I, § 9, no. 7, Proposition 10 and Chapter II, § 4, no. 3, Proposition 4).
\item
  Suppose in addition that Y is complete with respect to its uniformity; then A is also closed in X {[}observe that if \(x_{0} \in \overline{A}\) , then \(\mathrm{H}(x_{0})\) is relatively compact in Y by considering an ultrafilter on \(\mathrm{H}(x_{0})\) as the image of an ultrafilter on H{]}. In this case, if X is connected, then H is relatively compact in \(\mathcal{C}_{c}(\mathrm{X};\mathrm{Y})\) if and only if, for one point \(x_{0} \in X\) , the set \(\mathrm{H}(x_{0})\) is relatively compact in Y.
\item
  Take X to be the compact interval {[}0, 1{]} of R, and take Y to be the interval {]}0, 1{[} endowed with the uniformity induced by that of R. Give an example of an equicontinuous subset H of C (X; Y) such that the set A defined in a) is the interval {]}0, 1{]}.
\item
  Suppose that X = Y and that H consists of homeomorphisms of X onto itself. Show that if H is uniformly equicontinuous, then the set A defined in a) is both open and closed in X.
\end{enumerate}

\begin{enumerate}
\def\labelenumi{\arabic{enumi})}
\setcounter{enumi}{13}
\tightlist
\item
  Let \(\Gamma\) be an equicontinuous group of homeomorphisms of \(\mathbf{R}\) . Show that if a homeomorphism \(u \in \Gamma\) leaves at least one point of \(\mathbf{R}\) fixed, and if \(u\) is increasing, then \(u\) is the identity mapping (show that otherwise the group generated by \(u\) is not equicontinuous at a suitably chosen fixed point of \(u\) ). If \(u\) is decreasing, then \(u^2\) is the identity mapping.
\end{enumerate}

¶ 15) Let X be a compact metrizable space, let \(\Gamma\) be the group of all homeomorphisms of X, and let G be a subgroup of \(\Gamma\) which operates transitively on X. If H is the centralizer of G in \(\Gamma\) , show that H is equicontinuous. {[}Argue by reductio ad absurdum, supposing that H is not equicontinuous at some point \(a \in \mathbf{X}\) ; deduce the existence of a sequence of points \(x_n \in \mathbf{X}\) and a sequence of elements \(u_n \in \mathbf{H}\) such that

\[
\lim _ {n \rightarrow \infty} x _ {n} = a, \quad \lim _ {n \rightarrow \infty} u _ {n} (a) = b, \quad \lim _ {n \rightarrow \infty} u _ {n} (x _ {n}) = c,
\]

where \(b \neq c\) . Hence show that the sequence \((u_n)\) converges pointwise in \(X\) , but that no point of \(X\) is a point of uniform convergence of this sequence; this contradicts Exercise 9 of §1.{]}

\begin{enumerate}
\def\labelenumi{\arabic{enumi})}
\setcounter{enumi}{15}
\tightlist
\item
  Let X be a Hausdorff uniform space and let \(\Gamma\) be an equicontinuous group of homeomorphisms of X. The equivalence relation R defined by \(\Gamma\) on X is open (Chapter I, § 5, no. 2).
\end{enumerate}

\begin{enumerate}
\def\labelenumi{\alph{enumi})}
\item
  Show that if every orbit of \(\Gamma\) is closed in \(\mathbf{X}\) , then the orbit space \(\mathbf{X} / \Gamma\) is Hausdorff.
\item
  Show that if every orbit of \(\Gamma\) is compact, then the relation \(\mathbf{R}\) is closed (use Chapter I, § 5, no. 4, Proposition 10).
\item
  Give an example where X is compact but none of the orbits of \(\Gamma\) is closed in X (cf.~Chapter III, § 2, Exercise 29).
\end{enumerate}

¶ 17) Let X be a compact space and Γ a countable group of homeomorphisms of X. Suppose that the orbit space X/Γ is Hausdorff (and therefore compact).

\begin{enumerate}
\def\labelenumi{\alph{enumi})}
\item
  Show that, for each \(x \in \mathbf{X}\) , the orbit \(\Gamma(x)\) of \(x\) is finite (observe that if it were infinite it would have no isolated point, and use Baire's theorem (Chapter IX, § 5, no. 3, Theorem 1)).
\item
  For each \(x_0 \in \mathbf{X}\) , let \(\Delta(x_0)\) be the normal subgroup of \(\Gamma\) consisting of homeomorphisms which leave fixed all the points of the orbit \(\Gamma(x_0)\) . Show that if \(\Delta(x_0)\) is finitely generated, then \(\Gamma\) is equicontinuous at \(x_0\) . {[}Let \(f_i (1 \leqslant i \leqslant m)\) be the generators of \(\Delta(x_0)\) , and let \(g_k (1 \leqslant k \leqslant n)\) be representatives of each of the cosets of \(\Delta(x_0)\) in \(\Gamma\) other than \(\Delta(x_0)\) itself. Take a neighbourhood \(\mathbf{V}\) of \(x_0\) such that none of the sets \(f_i(\mathbf{V})\) , \(f_i^{-1}(\mathbf{V})\) meets any of the sets \(g_k(\mathbf{V})\) , and then use the fact that the equivalence relation defined by \(\Gamma\) is closed (Chapter I, § 10, no. 4, Proposition 8).{]}
\item
  Without making any hypothesis on \(\Delta(x_{0})\) , suppose that \(x_{0}\) has a countable fundamental system of connected neighbourhoods in X. Show that \(\Gamma\) is equicontinuous at \(x_{0}\) {[}same method as in b){]}. Deduce that if X is locally connected, then \(\Gamma\) is equicontinuous.
\item
  Take X to be the compact subspace of R consisting of the points 0, 1, 1/n and \(1 + 1/n\) (n an integer \(\geqslant 2\) ). Give an example of a countable group \(\Gamma\) of homeomorphisms of X, such that \(X/\Gamma\) is Hausdorff but \(\Gamma\) is not equicontinuous.
\end{enumerate}

¶ 18) Let G be a Hausdorff topological group operating continuously on a Hausdorff topological space X. G is said to be proper at a point \(x_{0} \in X\) if the orbit G. \(x_{0}\) is closed in X and if G operates properly on G. \(x_{0}\) (Chapter III, § 4, no. 1).

\begin{enumerate}
\def\labelenumi{\alph{enumi})}
\item
  Let G be a locally compact topological group operating continuously on a Hausdorff uniform space X. Suppose that the set of homeomorphisms \(x \rightarrow s.x\) of X, as s runs through G, is equicontinuous. Show that if G is proper at a point \(x_{0} \in X\) , then there is a neighbourhood V of \(x_{0}\) such that the set of elements \(s \in G\) such that \(s.V \cap V \neq \emptyset\) is relatively compact in G (cf.~Chapter III, § 4, no. 4, Proposition 7). Deduce that the set D of points of X at which G is proper is open in X and that G operates properly on D. If in addition X is locally compact, show that D is both open and closed in X (use no. 1, Proposition 1, Corollary 2).
\item
  In the real plane \(\mathbf{R}^2\) , let \(\mathbf{E}\) be the set consisting of the origin \((0,0)\) and the points \((0,2^{-n})\) as \(n\) runs through the integers \(\geqslant 0\) . For each \(n\in \mathbf{Z}\) such that \(n\neq 0\) let \(u_{n}\) be the restriction to \(\mathbf{E}\) of an affine linear mapping of \(\mathbf{R}^2\) into itself, such that \(u_{n}(0,0) = (n,0)\) , and \(u_{n}(0,2^{-|n|}) = (0,\theta^{n})\) where \(\theta\) is an irrational number such that \(0 <   \theta <  1\) . Let \(\mathbf{X}\) be the (locally compact) subspace of \(\mathbf{R}^2\) which is the union of \(\mathbf{E}\) and the sets \(u_{n}(\mathbf{E})\) for \(n\in \mathbf{Z},n\neq 0\) . Furthermore, let \(u_{0}\) denote the identity mapping of \(\mathbf{E}\) , and extend \(u_{n}(n\in \mathbf{Z})\) to the whole of \(\mathbf{X}\) by putting \(u_{n}(u_{m}(x)) = u_{n + m}(x)\) for all \(x\in \mathbf{E}\) and all \(m\in \mathbf{Z}\) . The set \(\mathbf{G}\) of mappings \(u_{n}\) is a group of homeomorphisms of \(\mathbf{X}\) which is given a discrete topology so that \(\mathbf{G}\) operates continuously on \(\mathbf{X}\) . Show that \(\mathbf{G}\) is proper at every point of \(\mathbf{X}\) , but is not equicontinuous and does not operate properly on \(\mathbf{X}\) .
\item
  Let X be the (not locally compact) subspace of \(R^{2}\) (identified with C) which is the union of the half-plane y \textgreater{} 0 and the origin. Define a homeomorphism u of X onto itself by putting \(u(0, 0) = (0, 0)\) and \(u(re^{i\omega}) = re^{i\omega'}\) where \(\omega' = \frac{1}{2}\omega\) if \(0 < \omega \leqslant \frac{1}{2}\pi\) , \(\omega' = \omega - \frac{1}{4}\pi\) for \(\frac{1}{2}\pi \leqslant \omega \leqslant \frac{3}{4}\pi\) , \(\omega' = 2\omega - \pi\) for \(\frac{3}{4}\pi \leqslant \omega < \pi\) (r \textless{} 0, \(0 < \omega < \pi\) ). Let G be the group of homeomorphisms of X generated by u, and endow G with the discrete topology. Show that G is equicontinuous on X and that the set of points of X at which G is proper is the half-plane y \textgreater{} 0.
\end{enumerate}

\begin{enumerate}
\def\labelenumi{\arabic{enumi})}
\setcounter{enumi}{18}
\tightlist
\item
  Let X be a Hausdorff uniform space and let G be a group of homeomorphisms of X onto itself.
\end{enumerate}

\begin{enumerate}
\def\labelenumi{\alph{enumi})}
\item
  Show that if \(\mathbf{G}\) is equicontinuous and discrete with respect to the topology of pointwise convergence, then \(\mathbf{G}\) is closed in \(\mathcal{F}_s(\mathbf{X};\mathbf{X})\) .
\item
  Show that if G, endowed with the discrete topology, is proper at one point of X at least (Exercise 18), then G is closed in \(\mathcal{F}_{s}(\mathbf{X};\mathbf{X})\) and the topology of pointwise convergence induces the discrete topology on G.
\item
  Let X be the topological sum of two spaces \(X_{1}, X_{2}\) , each homeomorphic to R, so that X can be identified with the product space \(R \times \{1, 2\}\) . Endow X with the product uniformity. For each pair \((m, n)\) of rational integers let \(u_{mn}\) denote the homeomorphism of X defined by \(u_{mn}(x_{1}) = x_{1} + m\alpha + n\beta\) , \(u_{mn}(x_{2}) = x_{2} + m\gamma + n\delta\) , where \(x_{1} \in X_{1}\) and \(x_{2} \in X_{2}\) , and \(\alpha, \beta, \gamma, \delta\) are four non-zero real numbers such that \(\alpha/\beta\) and \(\gamma/\delta\) are irrational and distinct. Show that the \(u_{mn}\) form a group of homeomorphisms of X which is equicontinuous and discrete with respect to the topology of pointwise convergence, but that G (with the discrete topology) is not proper at any point of X.
\end{enumerate}

\begin{enumerate}
\def\labelenumi{\arabic{enumi})}
\setcounter{enumi}{19}
\tightlist
\item
  Let G be a group of homeomorphisms of a Hausdorff uniform space X. Suppose that G is equicontinuous and that G, endowed with the discrete topology, operates properly on X. Suppose, furthermore, that the set of points \(x \in X\) , which are not fixed points of any homeomorphism \(u \in G\) other than the identity, is dense in X. Under these conditions, show that there is an open set F in X such that \(F \cap U(F) = \varnothing\) for all \(u \in G\) other than the identity, and such that the canonical image of F in the orbit space X/G is a dense open subset of X/G, homeomorphic to F (use Exercise 16 and Zorn's lemma).
\end{enumerate}

